\begin{abstract}
Whether a question is placed before or after an image changes nothing about a
prompt's content, yet it changes a vision-language model's accuracy. Prior work
by the present authors established this question-first paradox on three yes/no
benchmarks and traced it to read-out rather than perception: the question
demonstrably steers the image patches, but the answer position cannot read it
back, and echoing the question or the image repairs the damage. That work left
open how far the effect reaches, what the echo is actually made of, and what it
costs. We answer all three. Across fourteen benchmarks never previously tested,
question-first costs $7.0$ points on average and hurts on 12 of 13 image and
video datasets, and on a $24{,}449$-question presence benchmark it hurts on 7 of
7 categories, taking one from $86.3$ to $38.9$. It also has a limit we can name:
on RF20-VL object detection, where the output is bounding boxes rather than an
answer token, the sign reverses and question-first is $2.80$ mAP \emph{better},
winning on 17 of 20 datasets. The echo turns out to carry almost no information: downscaling the
echoed image to $\tfrac{1}{8}$ resolution, $\tfrac{1}{64}$ of the pixels, is
statistically indistinguishable from a full-resolution echo ($p = 1.00$) while
still beating question-last ($p = 0.005$), which cuts its price from $3{,}352$
extra tokens per query to $107$. Breadth also exposes a boundary: on video,
where the echo doubles an already long frame sequence, image echoing helps on
only 1 of 4 datasets while question echoing still recovers $3.3$ points. A
reflective prompt optimizer, given the same budget, gains only under the
ordering our diagnosis says has headroom and never reaches the question-last
baseline, so rewording does not substitute for fixing token order.
\end{abstract}
